%==============================================================================
% Chapter 1 — Mathematical Foundations for Machine Learning
% Curriculum: Phase 1, Chapter 1 of docs/AI_ML_Master_Checklist.md
%
% Figures for this chapter go in theory/figures/.
%==============================================================================
\chapter{Mathematical Foundations for Machine Learning}
\label{ch:math-foundations}

Before you can build the models, you need to speak the language they are written
in. This chapter fixes the notation used throughout the Compendium and works
through the pillars that every later chapter leans on. Three are structural ---
linear algebra, matrix calculus, and multivariate optimization --- and the rest
supply the vocabulary the book is stated in: numerical computing, probability,
statistical inference, and information theory.

Each result here reappears later in applied form. The singular value
decomposition returns as PCA in \cref{sec:unsupervised} and as low-rank
adaptation in \cref{sec:peft}; the chain rule returns as backpropagation in
\cref{sec:mlp}; Lagrangian duality returns as the SVM dual in \cref{sec:svm};
and the Kullback--Leibler divergence returns as both the cross-entropy loss and
the variational lower bound of \cref{sec:variational-inference}.

\section{Linear Algebra}
\label{sec:linear-algebra}

A vector $\vect{x} \in \R^{n}$ is a point in $n$-dimensional space; a matrix
$\mat{A} \in \R^{m \times n}$ is a linear map from $\R^{n}$ to $\R^{m}$. The
dot product of $\vect{x}, \vect{y} \in \R^{n}$,
%
\begin{equation}
    \inner*{\vect{x}, \vect{y}} = \vect{x}\T \vect{y}
        = \sum_{i=1}^{n} x_i y_i
        = \norm{\vect{x}}_2 \norm{\vect{y}}_2 \cos\theta ,
    \label{eq:dot-product}
\end{equation}
%
is the single most reused operation in machine learning. \Cref{eq:dot-product}
is why cosine similarity is the default metric for vector databases: normalise
both vectors and the dot product \emph{is} the cosine of the angle between them.

\subsection{The Singular Value Decomposition}

\begin{theorem}{Singular Value Decomposition}{svd}
    Every matrix $\mat{A} \in \R^{m \times n}$ admits a factorisation
    %
    \begin{equation}
        \mat{A} = \mat{U} \mat{\Sigma} \mat{V}\T ,
        \label{eq:svd}
    \end{equation}
    %
    where $\mat{U} \in \R^{m \times m}$ and $\mat{V} \in \R^{n \times n}$ are
    orthogonal ($\mat{U}\T\mat{U} = \mat{I}$, $\mat{V}\T\mat{V} = \mat{I}$) and
    $\mat{\Sigma} \in \R^{m \times n}$ is diagonal with non-negative entries
    $\sigma_1 \ge \sigma_2 \ge \dots \ge \sigma_r > 0$, $r = \rank(\mat{A})$.
\end{theorem}

The columns of $\mat{V}$ are the eigenvectors of $\mat{A}\T\mat{A}$, the columns
of $\mat{U}$ are the eigenvectors of $\mat{A}\mat{A}\T$, and the singular values
satisfy $\sigma_i = \sqrt{\lambda_i}$ for the shared non-zero eigenvalues
$\lambda_i$. Truncating \cref{eq:svd} to the largest $k$ singular values gives
the best rank-$k$ approximation of $\mat{A}$ in both the Frobenius and spectral
norms:
%
\begin{equation}
    \mat{A}_k = \sum_{i=1}^{k} \sigma_i \vect{u}_i \vect{v}_i\T ,
    \qquad
    \norm{\mat{A} - \mat{A}_k}_F^2 = \sum_{i=k+1}^{r} \sigma_i^2 .
    \label{eq:eckart-young}
\end{equation}

\begin{intuition}
    \Cref{eq:eckart-young} is the mathematical core of dimensionality reduction,
    image compression, and low-rank adaptation (LoRA). When a fine-tuning update
    is written as $\mat{W} = \mat{W}_0 + \mat{B}\mat{A}$ with
    $\mat{B} \in \R^{d \times k}$ and $\mat{A} \in \R^{k \times d}$, the claim
    being made is that the useful part of the update has small $k$ --- that its
    singular values decay fast.
\end{intuition}

\begin{example}[Rank-one structure]
    Let $\mat{A} = \vect{u}\vect{v}\T$ for non-zero $\vect{u}, \vect{v}$. Then
    $\rank(\mat{A}) = 1$, $\sigma_1 = \norm{\vect{u}}_2\norm{\vect{v}}_2$, and
    every remaining singular value is zero, so $\mat{A}_1 = \mat{A}$ exactly.
    No information is lost by storing $m + n$ numbers instead of $mn$.
\end{example}

\subsection{Verifying the Decomposition Numerically}

Every derivation in this book should be checkable in a few lines of code. The
convention throughout is that theory sections carry a runnable counterpart in
\texttt{projects/}.

\begin{lstlisting}[caption={Confirming the Eckart--Young bound from \cref{eq:eckart-young}.},label={lst:svd}]
import numpy as np


def truncated_svd(matrix: np.ndarray, rank: int) -> np.ndarray:
    """Return the best rank-`rank` approximation of `matrix`."""
    u, s, vt = np.linalg.svd(matrix, full_matrices=False)
    return (u[:, :rank] * s[:rank]) @ vt[:rank, :]


a = np.random.default_rng(0).standard_normal((8, 5))
_, sigma, _ = np.linalg.svd(a, full_matrices=False)

k = 2
error = np.linalg.norm(a - truncated_svd(a, k), ord="fro") ** 2
assert np.isclose(error, (sigma[k:] ** 2).sum())  # Eckart-Young
\end{lstlisting}

\section{Matrix Calculus}
\label{sec:calculus}

\begin{definition}{Gradient and Hessian}{grad-hess}
    For $f : \R^{n} \to \R$, the gradient collects the first partial
    derivatives and the Hessian collects the second:
    %
    \begin{equation*}
        \grad f(\vect{x}) =
        \begin{bmatrix} \pd{f}{x_1} \\[0.4em] \vdots \\[0.4em] \pd{f}{x_n} \end{bmatrix},
        \qquad
        \big[\Hess f(\vect{x})\big]_{ij} = \frac{\partial^2 f}{\partial x_i \partial x_j}.
    \end{equation*}
\end{definition}

The gradient points in the direction of steepest ascent, which is why gradient
\emph{descent} moves against it. The second-order Taylor expansion of $f$ about
$\vect{x}_0$,
%
\begin{equation}
    f(\vect{x}) \approx f(\vect{x}_0)
        + \grad f(\vect{x}_0)\T (\vect{x} - \vect{x}_0)
        + \tfrac{1}{2} (\vect{x} - \vect{x}_0)\T \Hess f(\vect{x}_0) (\vect{x} - \vect{x}_0),
    \label{eq:taylor}
\end{equation}
%
justifies both Newton's method and the second-order loss approximation that
XGBoost optimises at each boosting round.

\subsection{The Chain Rule}

Backpropagation is nothing more than the multivariate chain rule applied in
reverse topological order. For a composition
$\vect{z} = f_L \circ f_{L-1} \circ \dots \circ f_1(\vect{x})$, the Jacobian
factorises as
%
\begin{equation}
    \pd{\vect{z}}{\vect{x}}
        = \pd{f_L}{f_{L-1}} \, \pd{f_{L-1}}{f_{L-2}} \cdots \pd{f_1}{\vect{x}} .
    \label{eq:chain-rule}
\end{equation}

\begin{remark}
    Evaluating \cref{eq:chain-rule} right-to-left is forward-mode
    differentiation; left-to-right is reverse-mode. Reverse mode wins in deep
    learning because the loss is scalar, so the leftmost factor is a single row
    and every intermediate product stays a vector rather than a full matrix.
\end{remark}

\subsection{Layout Conventions for Vector and Matrix Derivatives}

% To write: numerator versus denominator layout, the Jacobian of a vector-valued
% map, and the identities for differentiating x^T A x and A x. Fixing the
% convention here prevents transpose errors in every later derivation.

\section{Multivariate Optimization}
\label{sec:optimization}

% Training is optimisation, so the vocabulary is set here once. Gradient descent
% appears in this section only as an optimisation scheme with convergence
% conditions; the algorithmic variants (momentum, RMSProp, Adam) belong to
% \cref{sec:optimizers}.

\subsection{Convexity}

\subsection{Stationary Points and Second-Order Conditions}

\subsection{Constrained Optimization: Lagrange Multipliers}

\subsection{The KKT Conditions}

\subsection{Gradient Descent as an Optimization Scheme}

\section{Numerical Computing and Stability}
\label{sec:numerical-stability}

% The gap between the mathematics above and the arithmetic a machine actually
% performs. Every bug in this section is one you will otherwise meet as a silent
% NaN eight chapters later: overflowing softmax, an ill-conditioned normal
% equation, a log-likelihood that underflows to zero.

\subsection{Floating-Point Representation and Machine Epsilon}

\subsection{Catastrophic Cancellation and Loss of Significance}

\subsection{The Log-Sum-Exp Trick}

\subsection{Condition Number: Why You Solve Rather Than Invert}

\section{Probability and Statistics}
\label{sec:probability}

Bayes' theorem inverts a conditional distribution:
%
\begin{equation}
    \Prob(\theta \given D)
        = \frac{\Prob(D \given \theta)\,\Prob(\theta)}
               {\int \Prob(D \given \theta)\,\Prob(\theta)\,\mathrm{d}\theta} .
    \label{eq:bayes}
\end{equation}

The numerator is cheap; the normalising integral in the denominator is the
reason MCMC and variational inference exist. Two estimators drop out of
\cref{eq:bayes} by discarding parts of it:
%
\begin{align}
    \hat{\theta}_{\mathrm{MLE}} &= \argmax_{\theta}\; \Prob(D \given \theta),
        \label{eq:mle} \\
    \hat{\theta}_{\mathrm{MAP}} &= \argmax_{\theta}\; \Prob(D \given \theta)\,\Prob(\theta).
        \label{eq:map}
\end{align}

Maximum likelihood ignores the prior entirely; MAP keeps it but reports only the
mode, throwing away the spread that makes the posterior useful. \Cref{tab:moments}
summarises the first two moments used throughout the book.

\begin{table}[H]
    \centering
    \caption{Moment definitions and their sample estimators.}
    \label{tab:moments}
    \begin{tabular}{@{}lll@{}}
        \toprule
        Quantity & Population & Sample estimator \\
        \midrule
        Mean       & $\mu = \E[X]$                    & $\bar{x} = \frac{1}{n}\sum_i x_i$ \\
        Variance   & $\Var(X) = \E[(X-\mu)^2]$        & $\frac{1}{n-1}\sum_i (x_i - \bar{x})^2$ \\
        Covariance & $\Cov(X,Y) = \E[(X-\mu_X)(Y-\mu_Y)]$ & $\frac{1}{n-1}\sum_i (x_i - \bar{x})(y_i - \bar{y})$ \\
        \bottomrule
    \end{tabular}
\end{table}

\section{Statistical Inference and Experiment Design}
\label{sec:statistical-inference}

% The frequentist counterpart to \cref{sec:probability}, and the section that
% answers the question every later chapter raises and cannot settle: "is model B
% actually better than model A, or did the test split flatter it?" A difference
% in validation score is an estimate with a standard error, not a fact.

\subsection{Sampling Distributions and the Central Limit Theorem}

\subsection{Confidence Intervals and the Bootstrap}

\subsection{Hypothesis Testing}

\subsection{P-Values and How They Are Misread}

\subsection{Multiple-Comparison Correction}

\subsection{Statistical Power, Effect Size and Sample Size}

\subsection{A/B Testing}

\section{Information Theory}
\label{sec:information-theory}

\begin{definition}{Entropy, Cross-Entropy and KL Divergence}{information}
    For discrete distributions $P$ and $Q$ over the same support $\mathcal{X}$,
    %
    \begin{align*}
        \Ent{P}    &= -\sum_{x \in \mathcal{X}} P(x)\log P(x), \\
        H(P, Q)    &= -\sum_{x \in \mathcal{X}} P(x)\log Q(x), \\
        \KL{P}{Q}  &= \sum_{x \in \mathcal{X}} P(x)\log\frac{P(x)}{Q(x)}
                    = H(P,Q) - \Ent{P} .
    \end{align*}
\end{definition}

Because $\Ent{P}$ does not depend on the model, minimising the KL divergence
between the empirical label distribution and the model's predictions is
identical to minimising cross-entropy. That equivalence is the reason
classifiers are trained with a cross-entropy loss at all, and it recurs when
deriving the evidence lower bound in Bayesian inference \autocite{bishop2006prml}.

\begin{intuition}
    $\KL{P}{Q}$ is the expected number of extra bits paid for encoding samples
    from $P$ with a code optimised for $Q$. It is non-negative, zero only when
    $P = Q$, and asymmetric --- $\KL{P}{Q} \neq \KL{Q}{P}$ in general. That
    asymmetry is a modelling choice, not a defect: forward KL spreads a fitted
    distribution across all modes, while reverse KL lets it collapse onto one.
\end{intuition}

\section{Notation and Figures}
\label{sec:notation}

\Cref{fig:placeholder} confirms that the figure pipeline resolves correctly:
images live in \texttt{figures/}, and \texttt{\textbackslash graphicspath} in
\texttt{preamble.tex} means the path in \texttt{\textbackslash includegraphics}
is a bare filename. Replace this placeholder with real diagrams as the chapter
is written out.

\begin{figure}[H]
    \centering
    \includegraphics[width=0.4\textwidth]{IMG_8232.jpg}
    \caption[Figure pipeline placeholder]{Placeholder image, included to verify
        that \texttt{\textbackslash graphicspath}, float placement, captions and
        cross-references all work end to end. Swap for a real diagram --- for
        this chapter, a geometric sketch of the SVD acting on the unit circle
        would be the natural first figure.}
    \label{fig:placeholder}
\end{figure}

The reference texts for this chapter are Strang \autocite{strang2019linalg} for
linear algebra, Cover and Thomas \autocite{cover2006information} for information
theory, and Goodfellow et al.\ \autocite{goodfellow2016deep} for the machine
learning framing of every section above.
